. wbopendata, indicator(SI.POV.DDAY;NY.GDP.PCAP.PP.KD) clear long
(using cached data: ind_si_pov_dday_all_en.csv, TTL 7d)
{\smallskip}
    {\bftt{Metadata for indicator}} SI.POV.DDAY
\HLI{80}
    {\bftt{Name}}: Poverty headcount ratio at $3.00 a day (2021 PPP) (\% of
    population)
\HLI{80}
    {\bftt{Collection}}: 2 World Development Indicators
\HLI{80}
    {\bftt{Description}}: Poverty headcount ratio at $3.00 a day is the percentage of
    the population living on less than $3.00 a day at 2021 purchasing power
    adjusted prices. As a result of revisions in PPP exchange rates, poverty
    rates for individual countries cannot be compared with poverty rates
    reported in earlier editions.
\HLI{80}
    {\bftt{Note}}: World Bank, Poverty and Inequality Platform. Data are based on
    primary household survey data obtained from government statistical
    agencies and World Bank country departments. Data for high-income
    economies are mostly from the Luxembourg Income Study database. For more
    information and methodology, please see http://pip.worldbank.org., World
    Bank (WB), uri: http://pip.worldbank.org, note: Data are based on
    primary household survey data obtained from government statistical
    agencies and World Bank country departments. Data for high-income
    economies are mostly from the Luxembourg Income Study database.
\HLI{80}
    {\bftt{Topic(}}{\sltt{s}}{\bftt{)}}: 11 Poverty ; 2 Aid Effectiveness ; 19 Climate Change
\HLI{80}
{\smallskip}
{\smallskip}
(using cached data: ind_ny_gdp_pcap_pp_kd_all_en.csv, TTL 7d)
{\smallskip}
    {\bftt{Metadata for indicator}} NY.GDP.PCAP.PP.KD
\HLI{80}
    {\bftt{Name}}: GDP per capita, PPP (constant 2021 international $)
\HLI{80}
    {\bftt{Collection}}: 2 World Development Indicators
\HLI{80}
    {\bftt{Description}}: This indicator provides values for gross domestic product
    (GDP) expressed in constant international dollars, converted by
    purchasing power parities (PPPs). PPPs account for the different price
    levels across countries and thus PPP-based comparisons of economic
    output are more appropriate for comparing the output of economies and
    the average material well-being of their inhabitants than exchange-rate
    based comparisons.
\HLI{80}
    {\bftt{Note}}: .
\HLI{80}
    {\bftt{Topic(}}{\sltt{s}}{\bftt{)}}:
\HLI{80}
{\smallskip}
{\smallskip}
{\smallskip}
